\section{Общие сведения}

\textbf{Вариант:} A-5.\\
\textbf{Группа:} A.\\
\textbf{Объём выборки:} \(n = 106\).\\
\textbf{Единицы измерения:} мс.\\
\textbf{Уровень значимости:} \(\alpha = 0.05\).\\
Данные представлены в виде четырёх столбцов \(X_1, X_2, X_3, X_4\). Для всех критериев используется двусторонняя альтернатива.

\section{Описательные статистики}

\begin{table}[H]
\centering
\caption{Числовые характеристики выборок}
\label{tab:stats}
\begin{tabular}{lcccc}
\toprule
Показатель & \(X_1\) & \(X_2\) & \(X_3\) & \(X_4\) \\
\midrule
Среднее (\(\bar{x}\)) & 67.200 & 69.102 & 74.725 & 6.779 \\
Медиана (\(m_e\)) & 67.630 & 68.700 & 73.695 & 4.835 \\
Несмещ. СКО (\(\hat{\sigma}\)) & 8.621 & 9.300 & 10.553 & 6.381 \\
Несмещ. дисперсия (\(\hat{\sigma}^2\)) & 74.329 & 86.489 & 111.360 & 40.720 \\
Минимум & 48.440 & 48.160 & 49.170 & 0.030 \\
Максимум & 86.930 & 90.790 & 96.910 & 38.200 \\
Q1 (25\%) & 61.333 & 61.710 & 68.958 & 1.858 \\
Q3 (75\%) & 73.040 & 75.595 & 81.738 & 9.820 \\
\bottomrule
\end{tabular}
\end{table}
